\section{Modul 06. Prosedur, fungsi, dan macro}

\subsection{Tujuan}

Setelah menyelesaikan modul ini, mahasiswa mampu:
\begin{enumerate}[nosep]
    \item mendefinisikan prosedur dengan \texttt{PROC} dan
          \texttt{ENDP};
    \item memanggil prosedur dengan \texttt{CALL} dan mengakhirinya
          dengan \texttt{RET};
    \item menyampaikan parameter ke prosedur melalui register;
    \item menjelaskan fungsi sebagai prosedur yang mengembalikan nilai
          di \texttt{AL} atau \texttt{AX};
    \item mendefinisikan dan memakai macro dengan \texttt{MACRO} dan
          \texttt{ENDM};
    \item merakit dan menjalankan program yang memakai prosedur,
          fungsi, dan macro.
\end{enumerate}

\subsection{Dasar singkat}

\subsubsection*{Prosedur dan pemanggilan}

Prosedur adalah sekelompok instruksi bernama yang dapat dipanggil
dari bagian lain program. Prosedur didefinisikan dengan \texttt{PROC}
dan diakhiri \texttt{ENDP}. Pemanggil memakai \texttt{CALL}, sedangkan
prosedur mengakhiri dirinya dengan \texttt{RET}. \texttt{CALL}
menyimpan alamat kembali di stack, dan \texttt{RET} mengambil alamat
itu. Pada modul ini \texttt{RET} ditulis tanpa angka.

Data disampaikan ke prosedur melalui register. Pemanggil mengisi
register sebelum \texttt{CALL}, dan prosedur membaca register itu.
Pada program pertama, \texttt{DX} berisi alamat string yang akan
dicetak. Prosedur yang sama dapat dipanggil berkali-kali dengan isi
\texttt{DX} yang berbeda, sehingga keluarannya berbeda pula.

\subsubsection*{Fungsi}

Fungsi adalah prosedur yang mengembalikan nilai di \texttt{AL} atau
\texttt{AX}. Selain menyiapkan register masukan, pemanggil juga harus
membaca register hasil setelah \texttt{CALL} kembali. Pada program
kedua, \texttt{AL} dipakai sebagai register masukan sekaligus
register hasil: nilai awal diletakkan di \texttt{AL}, fungsi
menambahkan 3, lalu hasilnya dibaca dari \texttt{AL}.

Karena prosedur dan pemanggil berbagi register, pemanggil perlu
mengetahui register mana yang dipakai untuk parameter dan mana yang
dipakai untuk hasil. Register lain yang tidak disepakati sebaiknya
tidak diandalkan tetap bernilai sama setelah pemanggilan.

\subsubsection*{Macro}

Macro didefinisikan dengan \texttt{MACRO} dan diakhiri \texttt{ENDM}.
Berbeda dari prosedur, macro tidak dipanggil dengan \texttt{CALL} dan
tidak diakhiri \texttt{RET}. Saat perakitan, assembler menyalin isi
macro ke setiap tempat macro itu ditulis. Karena itu, macro yang
dipakai berkali-kali memperbesar program, tetapi tidak memerlukan
lompatan saat program berjalan.

Macro dapat menerima parameter yang diganti secara tekstual. Pada
modul ini macro \texttt{CETAK\_KARAKTER} tidak memakai parameter;
karakter yang dicetak sudah disiapkan pada \texttt{DL}. Perbedaan
prosedur dan macro dirangkum pada Tabel~\ref{tab:m06-banding}.

\begin{table}[H]
    \centering
    \caption{Perbandingan prosedur dan macro.}
    \label{tab:m06-banding}
    \small
    \begin{tabular}{@{}lll@{}}
        \toprule
        \textbf{Aspek} & \textbf{Prosedur} & \textbf{Macro} \\
        \midrule
        Pemanggilan & \texttt{CALL} dan \texttt{RET} &
        nama macro ditulis langsung \\
        Waktu pengolahan & saat program berjalan &
        saat perakitan, isinya disalin \\
        Ukuran program & badan ditulis satu kali &
        disalin pada setiap pemakaian \\
        Parameter & melalui register &
        diganti secara tekstual \\
        \bottomrule
    \end{tabular}
\end{table}

Prosedur dipilih bila bagian tersebut panjang atau sering dipanggil,
sedangkan macro dipilih bila bagiannya pendek dan pemakaiannya
sedikit.

Rujukan singkat: Gunawan, Bagian 13; S'to, Bab XVI dan XVII; buku
ajar, \texttt{chapters/14} bagian prosedur dan macro.

\subsection{Alat dan bahan}

\begin{itemize}[nosep]
    \item Komputer dengan emu8086 terpasang.
    \item DOSBox dengan TASM, TLINK, dan DEBUG.
    \item Berkas \texttt{code/m06-prosedur.asm} dan
          \texttt{code/m06-fungsi-macro.asm}.
    \item Lembar jurnal praktikum.
\end{itemize}

\subsection{Contoh program}

Program pertama memakai satu prosedur pencetak. Prosedur
\texttt{CETAK} dipanggil dua kali dengan \texttt{DX} yang berbeda,
sehingga dua kalimat tercetak walaupun badan prosedur hanya ditulis
sekali.

\lstinputlisting[
    caption={Satu prosedur yang dipanggil dua kali.},
    label={lst:m06-prosedur}
]{code/m06-prosedur.asm}

Program kedua memakai satu fungsi dan satu macro. Fungsi
\texttt{TAMBAH\_TIGA} menerima nilai pada \texttt{AL} dan
mengembalikan hasilnya pada \texttt{AL}. Macro
\texttt{CETAK\_KARAKTER} mencetak karakter yang sudah berada pada
\texttt{DL}.

\lstinputlisting[
    caption={Fungsi pengembalian nilai di \texttt{AL} dan macro
    pencetak karakter.},
    label={lst:m06-fungsi-macro}
]{code/m06-fungsi-macro.asm}

\subsection{Langkah kerja}

\begin{enumerate}
    \item Siapkan folder kerja, lalu salin kedua berkas \texttt{.asm}
          ke dalamnya.
    \item Rakit dan tautkan program pertama. Jalankan kedua perintah
          berikut pada baris terpisah:
          \begin{quote}
          \texttt{tasm m06-prosedur.asm}\\
          \texttt{tlink m06-prosedur.obj}
          \end{quote}
          Jalankan hasilnya dan catat kedua kalimat yang tercetak.
    \item Jalankan program pertama dengan \emph{single step} pada
          emu8086. Catat isi \texttt{DX} sebelum setiap \texttt{CALL}
          dan tunjukkan bahwa eksekusi kembali ke instruksi sesudah
          \texttt{CALL} ketika \texttt{RET} dijalankan.
    \item Ubah isi \texttt{BARIS2}, rakit ulang, lalu jalankan.
          Pastikan kalimat kedua berubah tanpa mengubah prosedur.
    \item Rakit dan tautkan program kedua dengan perintah yang sama,
          lalu jalankan. Catat karakter yang tercetak.
    \item Ubah nilai awal \texttt{AL} dari 5 menjadi 2, rakit ulang,
          lalu jalankan. Catat keluarannya, lalu ubah penambahan di
          dalam fungsi dari 3 menjadi 4 dan catat kembali hasilnya.
\end{enumerate}

\subsection{Tugas}

\begin{enumerate}
    \item Jelaskan tugas \texttt{PROC}, \texttt{ENDP}, \texttt{CALL},
          dan \texttt{RET}.
    \item Jelaskan mengapa \texttt{TAMBAH\_TIGA} disebut fungsi.
          Sebutkan register masukan dan register hasilnya.
    \item Jelaskan perbedaan prosedur dan macro dari segi waktu
          pengolahan dan ukuran program.
    \item Jelaskan mengapa prosedur \texttt{CETAK} dapat mencetak dua
          kalimat berbeda tanpa diubah.
    \item Panggil macro \texttt{CETAK\_KARAKTER} dua kali dengan isi
          \texttt{DL} yang berbeda sehingga tercetak dua karakter
          berurutan.
\end{enumerate}

\subsection{Hasil yang dikumpulkan}

\begin{enumerate}
    \item Kedua berkas \texttt{.asm} beserta versi modifikasinya.
    \item Tangkapan layar keluaran kedua program.
    \item Tabel jejak pemanggilan berisi \texttt{DX} atau \texttt{AL}
          sebelum \texttt{CALL} dan hasil sesudah \texttt{RET}.
    \item Lembar jurnal berisi jawaban langkah kerja dan tugas.
\end{enumerate}
